% !TeX program = pdflatex
\documentclass[9pt,a4paper]{extarticle}
\usepackage[spanish,es-nodecimaldot,es-noshorthands]{babel}
\usepackage[T1]{fontenc}
\usepackage[utf8]{inputenc}
\usepackage{lmodern}
\renewcommand{\familydefault}{\sfdefault}
\usepackage[a4paper,left=0.92cm,right=0.92cm,top=0.72cm,bottom=0.78cm,headheight=9pt]{geometry}
\usepackage{graphicx}
\usepackage{xcolor}
\usepackage{tikz}
\usetikzlibrary{positioning,arrows.meta}
\usepackage{multicol}
\usepackage{tcolorbox}
\tcbuselibrary{skins}
\usepackage{booktabs}
\usepackage{tabularx}
\usepackage{array}
\usepackage{enumitem}
\usepackage{fancyhdr}
\usepackage{hyperref}
\usepackage{microtype}

\definecolor{cchenpurple}{HTML}{483888}
\definecolor{cchenlime}{HTML}{DDDC40}
\definecolor{textgray}{HTML}{333333}
\definecolor{mutedgray}{HTML}{6F6F6F}
\definecolor{lightgray}{HTML}{F3F3F5}
\definecolor{linegray}{HTML}{D9D9DF}
\definecolor{okgreen}{HTML}{2E7D32}
\definecolor{warnamber}{HTML}{F9A825}
\definecolor{badred}{HTML}{C62828}

\hypersetup{colorlinks=true, linkcolor=cchenpurple, urlcolor=cchenpurple,
  pdftitle={Brief ejecutivo - fuentes implementadas para extracción CCHEN-only},
  pdfauthor={CCHEN 360}}
\pagestyle{fancy}
\fancyhf{}
\fancyfoot[L]{\scriptsize\color{mutedgray}Observatorio CCHEN 360 -- extracción CCHEN-only}
\fancyfoot[R]{\scriptsize\color{mutedgray}\thepage}
\renewcommand{\headrulewidth}{0pt}
\renewcommand{\footrulewidth}{0.25pt}
\setlength{\parindent}{0pt}
\setlength{\parskip}{1.3pt}
\setlength{\columnsep}{0.42cm}
\hyphenpenalty=10000
\exhyphenpenalty=10000
\setlist[itemize]{leftmargin=*,topsep=1pt,itemsep=0.15pt,parsep=0pt}
\renewcommand{\arraystretch}{1.02}
\newcolumntype{Y}{>{\raggedright\arraybackslash}X}
\newcommand{\sectiontitle}[1]{\vspace{1pt}{\normalsize\bfseries\color{cchenpurple}#1}\par\vspace{0.2pt}{\color{cchenlime}\rule{0.25\linewidth}{1.1pt}}\vspace{0.8pt}}
\newcommand{\tightsection}[1]{\vspace{2pt}{\normalsize\bfseries\color{cchenpurple}#1}\par}
\newtcolorbox{keybox}{enhanced,colback=cchenpurple!5,colframe=cchenpurple,boxrule=0.5pt,arc=2pt,left=4pt,right=4pt,top=3pt,bottom=3pt}
\newtcolorbox{metricbox}[2]{enhanced,colback=#1!8,colframe=#1,boxrule=0.5pt,arc=2pt,left=3pt,right=3pt,top=3pt,bottom=3pt,title={#2},fonttitle=\bfseries\tiny\color{textgray},coltitle=textgray}
\newtcolorbox{calloutbox}[1]{enhanced,colback=lightgray,colframe=linegray,boxrule=0.4pt,arc=2pt,left=4pt,right=4pt,top=3pt,bottom=3pt,title={#1},fonttitle=\bfseries\scriptsize\color{cchenpurple},attach title to upper={\par\vspace{0.5pt}}}
\newcommand{\dotpattern}{\begin{tikzpicture}[remember picture,overlay]\foreach \x in {0.2,0.55,...,11.2} {\foreach \y in {-0.1,0.22,0.54,0.86} {\fill[cchenpurple!55] (\x,\y) circle (0.55pt); \fill[cchenlime!85] ({\x+0.16},{\y+0.13}) circle (0.48pt);}}\end{tikzpicture}}

\begin{document}
\thispagestyle{fancy}

\begin{minipage}[t]{0.34\textwidth}
  \includegraphics[height=0.92cm]{../../licitacion/assets/ministerio\_energia\_color.png}
\end{minipage}
\begin{minipage}[t]{0.36\textwidth}
  \vspace{0.08cm}\dotpattern
\end{minipage}
\begin{minipage}[t]{0.28\textwidth}
  \raggedleft\includegraphics[height=0.92cm]{../../licitacion/assets/cchen360\_logo\_color.png}
\end{minipage}

\vspace{0.12cm}
{\color{cchenpurple}\rule{\textwidth}{1.55pt}}
\vspace{0.05cm}

\begin{center}
{\fontsize{17}{18.5}\selectfont\bfseries\color{cchenpurple}Fuentes implementadas para extracción CCHEN-only}\\[-1pt]
{\fontsize{9.2}{10.6}\selectfont\color{textgray}Qué se descargó, dónde quedó guardado y cómo alimenta el Observatorio CCHEN 360}\\[-1pt]
{\fontsize{7.8}{9}\selectfont\color{mutedgray}Fecha de corte: 20 de mayo de 2026}
\end{center}
\vspace{0.02cm}

\begin{keybox}
\textbf{\color{cchenpurple}Mensaje central.} Se documentaron \textbf{22 fuentes implementadas}: \textbf{12 directas/API}, \textbf{6 derivadas por semillas}, \textbf{2 runtime/local}, \textbf{1 bloqueada por credencial} y \textbf{1 match a revisar}. No se descarga el universo completo de cada fuente: se extraen registros con señal CCHEN por DOI, autor, afiliación, ROR/ORCID, alias institucional o semillas temáticas justificadas.
\end{keybox}

\vspace{0.02cm}
\begin{minipage}{0.24\textwidth}
\begin{metricbox}{okgreen}{Fuentes documentadas}
{\LARGE\bfseries 22}\\[-2pt]
{\tiny implementadas o registradas}
\end{metricbox}
\end{minipage}\hfill
\begin{minipage}{0.24\textwidth}
\begin{metricbox}{cchenpurple}{Directas/API}
{\LARGE\bfseries 12}\\[-2pt]
{\tiny listas para mostrar primero}
\end{metricbox}
\end{minipage}\hfill
\begin{minipage}{0.24\textwidth}
\begin{metricbox}{warnamber}{Con observación}
{\LARGE\bfseries 8}\\[-2pt]
{\tiny semillas, token o match}
\end{metricbox}
\end{minipage}\hfill
\begin{minipage}{0.24\textwidth}
\begin{metricbox}{badred}{Bloqueo crítico}
{\LARGE\bfseries 1}\\[-2pt]
{\tiny PatentsView requiere API key}
\end{metricbox}
\end{minipage}

\begin{minipage}[t]{0.48\textwidth}

\sectiontitle{1. Qué se descargó}

\begin{tabularx}{\linewidth}{@{}>{\raggedright\arraybackslash}p{0.29\linewidth}Y@{}}
\toprule
\textbf{Tipología} & \textbf{Datos descargados y uso} \\
\midrule
Publicaciones y citas & OpenAlex, CrossRef, Semantic Scholar, PubMed, Europe PMC, INSPIRE y arXiv consolidan DOI, autores, afiliaciones, conceptos, citas, abstracts y preprints CCHEN. \\
Investigadores y acceso & ORCID mantiene perfiles de investigadores; Unpaywall marca acceso abierto por DOI. \\
Datasets y outputs & DataCite, OpenAIRE y Zenodo inventarian datasets, presentaciones, archivos y outputs asociados a CCHEN. \\
Bio/Farma nuclear & PubChem/radiofarmacia descarga semillas, compuestos/radionúclidos y literatura técnica; ClinVar, GenBank, GEO, NIH y SRA quedan como evidencia derivada. \\
Patentes & INAPI queda como base local; PatentsView/USPTO está registrado como fuente API, pero bloqueado hasta contar con credencial gratuita. \\
Institucional y vigilancia & Datos.gob.cl, IAEA INIS y noticias CCHEN/nuclear agregan convenios, acuerdos, registros nucleares y señales de vigilancia. \\
\bottomrule
\end{tabularx}

\begin{calloutbox}{Volumen de evidencia}
Conteos principales: OpenAlex (949; 21348; 877; 8499); Crossref (1434); Semantic Scholar (1057); OpenAIRE (878); PubChem (875); Unpaywall (780); arXiv (407; 1); IAEA (INIS) (110); Zenodo (68); ORCID (65). Los valores múltiples son artefactos distintos y no se suman como una sola tabla.
\end{calloutbox}

\sectiontitle{2. Dónde quedaron guardados}

\begin{itemize}
  \item \textbf{Publicaciones:} \texttt{Data/Publications/}.
  \item \textbf{Outputs:} \texttt{Data/ResearchOutputs/}.
  \item \textbf{Gobernanza/radiofarmacia:} \texttt{Data/Gobernanza/}.
  \item \textbf{Vigilancia/patentes:} \texttt{Data/Vigilancia/} y \texttt{Data/Patents/}.
\end{itemize}

\end{minipage}\hfill
\begin{minipage}[t]{0.48\textwidth}

\sectiontitle{3. Estado y frecuencia de descarga}

\begin{tabularx}{\linewidth}{@{}>{\raggedright\arraybackslash}p{0.31\linewidth}Y@{}}
\toprule
\textbf{Frecuencia} & \textbf{Fuentes} \\
\midrule
Semanal & arXiv monitor, IAEA INIS y noticias CCHEN/nuclear. \\
Trimestral & OpenAlex, CrossRef, Semantic Scholar y radiofarmacia/semillas. \\
Semestral & ORCID, PubMed, Europe PMC, INSPIRE, DataCite, OpenAIRE, Unpaywall, Zenodo, datos institucionales y patentes. \\
\bottomrule
\end{tabularx}

\sectiontitle{4. Uso para el observatorio}

\begin{itemize}
  \item \textbf{Producción científica:} alimentar indicadores de publicaciones, autores, colaboración, citas, temas y acceso abierto.
  \item \textbf{Radiofarmacia y medicina nuclear:} vigilar compuestos, radionúclidos, literatura técnica y oportunidades temáticas.
  \item \textbf{Vigilancia institucional:} monitorear señales nucleares, prensa CCHEN/nuclear y oportunidades de posicionamiento.
  \item \textbf{Gobernanza de datos:} evidenciar qué fuentes ya existen, qué producen y qué requiere monitoreo.
\end{itemize}

\begin{calloutbox}{Observaciones críticas}
\textbf{PatentsView/USPTO} requiere \texttt{PATENTSVIEW\_API\_KEY} y no reemplaza INAPI local. \textbf{Google Finance / News monitor} debe revisarse: el runtime implementado monitorea noticias CCHEN/nuclear, no datos financieros. Las fuentes Life Sciences derivadas deben presentarse como evidencia temática hasta implementar extractores directos.
\end{calloutbox}

\sectiontitle{5. Entregables}
\textbf{CSV operativo}, \textbf{briefs PDF} y \textbf{gráficos de estado}. Cada respaldo incluye tipología, uso, frecuencia, estado, comentario y decisión.

\end{minipage}

\vspace{0.08cm}
\begin{tcolorbox}[enhanced,colback=white,colframe=cchenpurple!55,boxrule=0.45pt,arc=2pt,left=5pt,right=5pt,top=4pt,bottom=4pt]
{\bfseries\color{cchenpurple}Método replicable de extracción:}\hfill
\begin{tikzpicture}[baseline=-0.55ex, every node/.style={font=\scriptsize, align=center}]
\node[draw=cchenpurple,fill=cchenpurple!8,rounded corners=2pt,inner xsep=5pt,inner ysep=3pt] (a) {Fuente API\\o archivo local};
\node[draw=okgreen,fill=okgreen!8,rounded corners=2pt,inner xsep=5pt,inner ysep=3pt,right=0.32cm of a] (b) {Filtro\\CCHEN-only};
\node[draw=cchenpurple,fill=cchenpurple!8,rounded corners=2pt,inner xsep=5pt,inner ysep=3pt,right=0.32cm of b] (c) {CSV local\\y runtime};
\node[draw=warnamber,fill=warnamber!10,rounded corners=2pt,inner xsep=5pt,inner ysep=3pt,right=0.32cm of c] (d) {Calidad\\y brechas};
\node[draw=okgreen,fill=okgreen!8,rounded corners=2pt,inner xsep=5pt,inner ysep=3pt,right=0.32cm of d] (e) {Indicador\\Observatorio};
\draw[-{Latex[length=2mm]},cchenpurple] (a) -- (b);
\draw[-{Latex[length=2mm]},cchenpurple] (b) -- (c);
\draw[-{Latex[length=2mm]},cchenpurple] (c) -- (d);
\draw[-{Latex[length=2mm]},cchenpurple] (d) -- (e);
\end{tikzpicture}
\end{tcolorbox}
\end{document}
